\section{Discussion}
\label{sec:discussion}

\subsection{The Em Dash as Fine-Tuning Signature}

The most striking empirical finding is not that LLMs produce em dashes, but that the pattern of production---frequency, suppressibility, and provider variation---constitutes a signature of the post-training process. While we use ``RLHF'' as shorthand throughout, we note that providers differ on many axes beyond reward modeling: training data curation, supervised fine-tuning data, constitutional AI methods, tokenizer design, and system-prompt defaults all contribute to the final output distribution. The em dash signature reflects the aggregate effect of these decisions, not RLHF in isolation. GPT-4.1 retains 9.10 em dashes per 1,000 words even under explicit suppression; DeepSeek~V3 retains 5.41; Claude Opus~4.6 drops to 0.19; Gemini~2.5~Pro drops to zero; Llama produces none; and OpenAI's newest model (GPT-5.4) has been reduced to 0.29, suggesting the mechanism is already being addressed across model generations. If the em dash were simply a training data artifact, we would expect it to appear uniformly across models trained on similar corpora. Instead, it appears selectively, at rates determined by post-training decisions. This is consistent with the genealogy's architecture: training data installs the latent tendency; RLHF determines the amplitude.

This has implications beyond the em dash itself. If fine-tuning procedures leave characteristic marks in the prose they produce, then stylistic analysis of model outputs can function as a form of \textit{model attribution}---identifying not only whether text was machine-generated, but which provider's fine-tuning produced it.

\subsection{Alternative Hypotheses}

Goedecke~\cite{goedecke2025emdash} has proposed an alternative account: that em dash frequency in LLM output reflects the prevalence of em dashes in digitized 19th-century literature included in training corpora, an era when em dash usage peaked. This account is not incompatible with ours. The markdown-training genealogy and the historical-literature hypothesis identify different layers of the training distribution that may both contribute to the latent tendency that post-training then amplifies. Our account adds the structural mechanism---why the em dash specifically survives formatting suppression---which the historical-literature account does not address.

A second alternative is that the em dash functions as a prestige marker in written English---associated with educated, literary, and formally articulate prose---and that RLHF amplifies it not because of any structural orientation from markdown but simply because evaluators reward prose that ``sounds smart.'' On this account, the em dash is selected for its stylistic connotations rather than its structural ancestry. This hypothesis is also not incompatible with ours: the em dash's prestige associations and its structural role in training data may both contribute to its selection during fine-tuning. However, the prestige-marker account alone does not explain why the em dash specifically resists formatting suppression while other prestige markers (e.g., semicolons, subordinate clause structures) do not show comparable resistance.

\subsection{Implications for AI Writing Detection}

Current approaches to AI-generated text detection focus primarily on statistical properties of token distributions~\cite{mitchell2023detectgpt} or deliberately embedded watermarks~\cite{kirchenbauer2023watermark}. The ``last fingerprint'' framework suggests a complementary signal: \textit{structural-register markers}---prose-level artifacts that betray an underlying structural orientation inherited from training data and shaped by fine-tuning.

The suppression test results demonstrate that these markers are robust to explicit instructions to avoid formatting. The em dash survives scrubbing because the scrubbing instruction does not recognize it as formatting. Furthermore, the finding that different providers produce distinctive em dash signatures suggests that structural-register analysis could contribute to model attribution as well as binary human-vs-machine classification.

\subsection{Future Work}

Two experiments would further extend these findings. First, a \textbf{base-vs-instruct comparison in an amplifying family}: our current comparison (Table~\ref{tab:base}) uses Llama, where fine-tuning suppresses em dashes; a comparison within a family where fine-tuning amplifies them would more directly isolate the amplification step. This requires access to base weights from OpenAI, Anthropic, or DeepSeek, which are not currently public. Second, a \textbf{training corpora analysis}: quantifying the frequency with which dashes (\texttt{---}, \texttt{--}, em dash characters) appear at structural boundaries versus inline positions in major training datasets would test the genealogy's prediction that structural-boundary dashes dominate, establishing the statistical basis for the model's association between dashes and structural transitions.

\subsection{Limitations}

\begin{enumerate}[leftmargin=*]
\item \textbf{Limited base model comparison.} We include a base-vs-instruct comparison for Llama~3.1~8B (Table~\ref{tab:base}), confirming that a latent tendency exists pre-RLHF. However, this comparison is available only for the Llama family, where RLHF suppresses rather than amplifies em dashes. A base-vs-instruct comparison within a family where RLHF amplifies the artifact (e.g., an OpenAI or Anthropic model, if base weights were available) would more directly isolate the amplification step.

\item \textbf{The Pandoc rendering claim is supporting, not primary.} Some markdown processors (notably Pandoc with smart typography~\cite{pandoc2024}) convert \texttt{---} to em dashes inline. However, the most widely used processors---including CommonMark~\cite{commonmark2021} and GitHub Flavored Markdown~\cite{gfm2024}---do not. The core argument does not depend on any specific rendering pipeline.

\item \textbf{Human baseline variance.} Em dash frequency in human prose ranges from 0.33 to 17.12 per 1,000 words in our sample, a 50$\times$ range. Claims about ``overuse'' relative to human norms must be qualified by genre and author. The paper's central findings---suppression asymmetry and provider variation---hold independently of the absolute human baseline.

\item \textbf{Sample size and prompt sensitivity.} Individual samples show high variance. Larger-scale replication with more diverse prompts would strengthen confidence in the frequency estimates, particularly for models near the human baseline.

\item \textbf{Falsifiability.} The complementarity argument (Section~\ref{sec:genealogy}) is difficult to test directly: if a model shows elevated em dashes, the genealogy attributes this to markdown-trained orientation amplified by fine-tuning; if it does not, the genealogy attributes this to fine-tuning suppressing the tendency. A stronger test would require a model family trained on corpora with minimal markdown content (e.g., predominantly book-text or speech transcripts) where the genealogy predicts low baseline em dash rates even before fine-tuning. Such corpora conditions are not currently available for controlled comparison.
\end{enumerate}
